\documentclass[11pt,a4paper]{article}
\usepackage[T1]{fontenc}
\usepackage[utf8]{inputenc}
\usepackage{lmodern}
\usepackage[margin=27mm]{geometry}
\usepackage{amsmath,amssymb,amsthm,bm,booktabs}
\usepackage{microtype}
\usepackage[numbers,sort&compress]{natbib}
\usepackage[hidelinks]{hyperref}
\newcommand{\Tr}{\operatorname{Tr}}
\newcommand{\diag}{\operatorname{diag}}
\newcommand{\vecop}{\operatorname{vec}}
\newcommand{\dist}{\operatorname{dist}_F}
\newcommand{\norm}[1]{\lVert #1\rVert_F}
\newcommand{\ip}[2]{\langle #1,#2\rangle}
\newcommand{\Mq}{\mathcal M_q}
\newtheorem{theorem}{Theorem}
\newtheorem{proposition}[theorem]{Proposition}
\newtheorem{lemma}[theorem]{Lemma}
\title{Local stability of a sharp \(q\)-deformed commutator response\\
for conditional quantum operations}
\author{Xiaojian Zeng\\
\small PBC School of Finance, Tsinghua University, Beijing 100190, China\\
\small \texttt{xiaojianzeng@foxmail.com}}
\date{}
\begin{document}
\maketitle
\begin{abstract}
The real overlap of an ordinary commutator with a \(q\)-deformed
commutator has a known sharp upper bound for two-dimensional matrices
and for normal matrices in arbitrary dimension. We study its optimized
response near an embedded two-level saturation witness. For every fixed
finite dimension \(n\geq3\) and every real \(q\notin\{-1,1\}\), we prove
that the response has a Morse--Bott local maximum at the known sharp
constant. Its local equality set is a smooth real manifold of dimension
\(4n-2\), characterized by one complex relation on a reducing
two-dimensional block. We calculate the complete normal Hessian spectrum,
including perturbations that couple the active levels to additional levels,
and obtain a two-sided quadratic deficit in the Frobenius distance to
the equality manifold. The resulting neighborhood includes nonnormal,
full-rank matrices, while the second operator remains unrestricted.
An auxiliary determinant identity characterizes all nonzero
two-dimensional equality matrices at \(q=1/2\).
After contraction scaling, the response is a success-weighted observable
of coherently ordered conditional operations on a maximally entangled
input. The theorem supplies a local sharp envelope and exact normal
susceptibilities for that observable. No explicit neighborhood radius,
uniform limit, or resolution of the unrestricted global conjecture is
claimed.
\end{abstract}
\noindent\textbf{Keywords:} conditional quantum operations;
deformed commutator; Frobenius norm; equality manifold;
local stability; operator interference.

\section{Introduction}
Operator products in different orders give distinct conditional
amplitudes. Their squared norms and real overlap can be combined into
a weighted order response. A useful matrix functional for this purpose is
\begin{equation}
 f(A,B;q)=\operatorname{Re}\Tr\bigl((BA-AB)^\dagger(BA-qAB)\bigr),
 \qquad q\in\mathbb R.
 \label{eq:f}
\end{equation}
Here \(A,B\in M_n(\mathbb C)\), the inner product is
\(\ip{X}{Y}=\Tr(X^\dagger Y)\), and
\(\norm{X}^2=\Tr(X^\dagger X)\) is the unnormalized Frobenius norm.
The parameter \(q\) is a real response weight.

Chru\'sci\'nski, Kimura, Ohno and Singal introduced this functional and
proved the sharp estimate
\begin{equation}
 f(A,B;q)\leq c(q)\norm{A}^2\norm{B}^2,\qquad
 c(q)=\frac{1+q+\sqrt{2(1+q^2)}}2,
 \label{eq:global}
\end{equation}
when \(n=2\), and when \(A\) is normal in any dimension
\citep[Theorem~1]{ckos2022response}. Their unrestricted all-dimensional
statement is a conjecture in that source. We use its author-preprint
version for the equations, equality witness and two-dimensional proof.
The ordinary B\"ottcher--Wenzel point \(q=1\) and the point \(q=-1\)
are already known cases.

The question addressed here is local: does the sharp witness remain
a maximum under arbitrary small matrix perturbations, including
nonnormal perturbations and couplings to extra levels?
We answer this for every fixed finite \(n\geq3\) and every real
\(q\notin\{-1,1\}\). The zero directions of the Hessian form a true
curved equality manifold, rather than an unresolved higher-order
instability. The remaining directions have an explicit strictly
negative spectrum. Together with a simple spectral branch, this
controls the full local remainder and proves a sharp local bound
for all complex \(B\).

This functional differs from the squared norm of a
\(q\)-deformed commutator studied in
\citep{ckos2022norm,nobori2026}. It also differs from the averaged
relaxation functional used in open quantum dynamics. The universal
maximal relaxation-rate bound has been proved in
\citep{muratore2024}; the present result supplies no new relaxation-rate
bound. Its physical content is instead the local envelope and normal
susceptibility of a success-weighted conditional-operation response.
The unrestricted inequality \eqref{eq:global} is not proved here.

\section{Response operator and local stability theorem}
Let column vectorization identify \(M_n(\mathbb C)\) with
\(\mathbb C^{n^2}\). Right and left multiplication by \(A\) are
represented by
\begin{equation}
 P_A=A^{\mathsf T}\otimes I_n,\qquad Q_A=I_n\otimes A.
\end{equation}
Expanding \eqref{eq:f} gives the Hermitian response operator
\begin{equation}
 H_q(A)=P_A^\dagger P_A+qQ_A^\dagger Q_A
       -\frac{1+q}{2}(P_A^\dagger Q_A+Q_A^\dagger P_A).
 \label{eq:H}
\end{equation}
Thus
\begin{equation}
 \begin{aligned}
 f(A,B;q)&=\vecop(B)^\dagger H_q(A)\vecop(B),\\
 g_q(A)&=\frac{\lambda_{\max}(H_q(A))}{\norm{A}^2}
       =\max_{B\neq0}\frac{f(A,B;q)}{\norm{A}^2\norm{B}^2}.
 \end{aligned}
 \label{eq:g}
\end{equation}
The function \(g_q\) is defined for \(A\neq0\), invariant under unitary
conjugation, and unchanged by multiplying \(A\) by a nonzero complex
scalar. No restriction on \(B\) is imposed by this optimization.

Fix \(q\notin\{-1,1\}\), and abbreviate
\begin{equation}
 \begin{gathered}
 R=\sqrt{2(1+q^2)},\quad c=c(q),\quad
 u=\sqrt{c-1},\quad b=\operatorname{sgn}(1+q)\sqrt{c-q},\\
 A_q=\diag(u,-b,0,\ldots,0).
 \end{gathered}
 \label{eq:params}
\end{equation}
Useful exact relations are
\begin{equation}
 \begin{gathered}
 u^2+b^2=R,\qquad u^2-b^2=q-1,\qquad
 2ub=1+q,\qquad c>\max(1,q).
 \end{gathered}
 \label{eq:relations}
\end{equation}
In particular \(u,b,u-b,u+b\) are nonzero on this parameter domain.

For a complex \(2\times2\) matrix \(C\), set
\begin{equation}
 \Phi_q(C)=\det C+\frac{ub}{(u-b)^2}(\Tr C)^2.
 \label{eq:Phi}
\end{equation}
Near \(A_q\), define \(\Mq\) to be the local family of matrices
supported on a reducing complex two-plane, with active block \(C\)
near \(\diag(u,-b)\) and \(\Phi_q(C)=0\).
Being supported means that in an orthonormal decomposition into this
plane and its complement the matrix is \(C\oplus0\).
Distance to \(\Mq\) below is distance to this local manifold in
a sufficiently small tubular neighborhood.
We use the real Frobenius metric
\(\operatorname{Re}\Tr(D^\dagger E)\) on matrix variations.

\begin{theorem}[Local sharp envelope and normal spectrum]
\label{thm:local}
For every fixed finite \(n\geq3\) and every real
\(q\notin\{-1,1\}\), the function \(g_q\) is real analytic near
\(A_q\) and has a Morse--Bott local maximum \(c\).
Its local maximum set and local critical set are exactly \(\Mq\),
a smooth embedded real manifold of dimension \(4n-2\).
At \(A_q\), the nonzero eigenvalues of its real Hessian are
\begin{equation}
 \begin{array}{c|c}
 \text{eigenvalue}&\text{multiplicity}\\ \hline
 -2&2\\
 -a_1=-\dfrac{4c(c-q)}{(2c-q)R}&2(n-2)\\[2mm]
 -a_2=-\dfrac{4c(c-1)}{(2c-1)R}&2(n-2)\\[2mm]
 -a_3=-\dfrac{2c}{R}&2(n-2)^2
 \end{array}
 \label{eq:spectrum}
\end{equation}
All four magnitudes are positive. With
\(\mu=\min(2,a_1,a_2,a_3)\) and
\(\Lambda=\max(2,a_1,a_2,a_3)\), there is a sufficiently small,
\(n,q\)-dependent neighborhood \(U\) of \(A_q\) on which
\begin{equation}
 \frac{\mu}{8}\dist(A,\Mq)^2
 \leq c-g_q(A)\leq
 \Lambda\dist(A,\Mq)^2.
 \label{eq:deficit}
\end{equation}
Consequently, for every \(A\in U\) and every \(B\in M_n(\mathbb C)\),
\begin{equation}
 f(A,B;q)\leq
 \left[c-\frac{\mu}{8}\dist(A,\Mq)^2\right]
 \norm{A}^2\norm{B}^2.
 \label{eq:localf}
\end{equation}
\end{theorem}

The theorem asserts existence of a neighborhood, not an explicit
radius. It contains full-rank and fully nonnormal perturbations of
\(A_q\), including arbitrary couplings to the additional levels.
There is no claim of uniformity as \(n\) grows or \(q\) approaches
an excluded point. The distance estimate is expressed at the
unnormalized witness \eqref{eq:params}; its constants refer to that
choice of Frobenius coordinates.

\section{Equality geometry and the largest spectral branch}
\begin{lemma}[Embedded equality family]
\label{lem:family}
The family \(\Mq\) is smooth of real dimension \(4n-2\).
On a neighborhood of \(A_q\) within \(\Mq\), one has
\(g_q=c\), the largest eigenvalue of \(H_q\) is simple,
and the full gradient of \(g_q\) vanishes.
\end{lemma}
\begin{proof}
The two-dimensional saturation witness of
\citep{ckos2022response} admits the Schur family
\begin{equation}
 C(w)=\begin{pmatrix}u&w\\0&-b\end{pmatrix},\qquad
 B(w)=\begin{pmatrix}-bw/R&1\\0&uw/R\end{pmatrix},
 \qquad w\in\mathbb R.
 \label{eq:pair}
\end{equation}
Substitution into \eqref{eq:H}, using \eqref{eq:relations}, gives
\begin{equation}
 H_q(C(w))\vecop(B(w))
       =c(R+w^2)\vecop(B(w)),\qquad \Phi_q(C(w))=0.
 \label{eq:paireigen}
\end{equation}
The known two-dimensional bound makes this the largest eigenvalue.
This verification uses the signed \(b\) and therefore includes
\(q<-1\).

The polynomial \(\Phi_q\) depends only on trace and determinant.
For its zeros near \(\diag(u,-b)\), the two unordered eigenvalues
are \(\alpha u,-\alpha b\), with \(\alpha\neq0\).
Indeed their sum is \(\alpha(u-b)\), and their product is
\(-\alpha^2ub\); conversely these roots give \(\Phi_q=0\).
Unitary Schur reduction, a diagonal unitary and complex scaling
then reduce such a block to \eqref{eq:pair} with \(w\geq0\).
Thus every block in the local \(\Phi_q\) family attains \(g_q=c\).

At the base, differentiating \(\Phi_q=0\) gives
\begin{equation}
 \delta C_{22}=-\frac bu\,\delta C_{11},
 \qquad \delta C_{12},\delta C_{21}\in\mathbb C.
 \label{eq:tangent}
\end{equation}
The complex gradient is nonzero, so the active equality set
has real dimension six. The Grassmannian of complex two-planes
in \(\mathbb C^n\) contributes \(4(n-2)\) dimensions.
The active block remains invertible nearby. Its supporting plane
is therefore uniquely recovered as the orthogonal complement of
the kernel of the full matrix. This proves local embeddedness
and the stated dimension.

At \(A_q\), the matrix units diagonalize \(H_q\), and the top
eigenvector is \(E_{12}\), with eigenvalue \(cR\).
The gap to \(E_{21}\) is \((1-q)^2>0\).
The other nonzero channels have eigenvalues among
\(u^2,b^2,qu^2,qb^2\), all strictly below \(cR\).
The remaining channels have eigenvalue zero.
This establishes simplicity, which persists near the base.

For completeness, separation also holds on nearby embedded
equality blocks. The two off-plane forms are right multiplication
with response \(\norm{BC}^2\) and left multiplication with
response \(q\norm{CB}^2\).
The first has top eigenvalue at most \(\norm{C}^2\).
For \(q\geq0\), the second has top eigenvalue at most
\(q\norm{C}^2\); for \(q<0\), it is negative semidefinite
and its top eigenvalue is at most zero.
Both are strictly below \(c\norm{C}^2\).
The active largest branch remains the simple one found above.

Active-block derivatives of \(g_q\) vanish because the known
two-dimensional bound is attained there. For off-plane and
bottom-block variations, the first variations of \(BA\) and
\(AB\), evaluated at an active maximizing \(B\), are orthogonal
to the active products. The first variation of \(\norm{A}^2\)
also vanishes in these directions. The simple-eigenvalue
derivative therefore makes the full gradient zero.
Unitary invariance extends this argument to every nearby
supporting plane.
\end{proof}

\section{Complete normal Hessian}
We give the calculation in real and imaginary matrix-entry
coordinates; \(E_{ij}\) and \(iE_{ij}\) form a real orthonormal basis.
Write \(\lambda(A)=\lambda_{\max}(H_q(A))\) on its simple branch,
and \(e=\vecop(E_{12})\).
For an entry direction \(D_i\), let \(P_i,Q_i\) denote its
right and left multiplication operators, and set
\(L_i=P_i-Q_i\), \(M_i=P_i-qQ_i\).
For the corresponding derivative \(H_i\) at \(A_q\),
simple-eigenvalue differentiation gives
\begin{align}
 \lambda_{ij}
 &=\operatorname{Re}\bigl(
   \ip{L_i e}{M_j e}+\ip{L_j e}{M_i e}\bigr)
   +2\operatorname{Re}\sum_{r\neq12}
       \frac{(H_i)_{12,r}(H_j)_{r,12}}
            {cR-(H_q(A_q))_{r,r}}, \label{eq:lambdahess}\\
 (g_q)_{ij}
 &=\frac{\lambda_{ij}-2c\delta_{ij}}R .
 \label{eq:quotienthess}
\end{align}
The last formula uses stationarity at \(A_q\),
\(\lambda=c\norm{A_q}^2\), and the Euclidean Hessian
of the squared Frobenius norm.

On the active real diagonal coordinates the Hessian block is
\begin{equation}
 -\frac2R\begin{pmatrix}b^2&ub\\ub&u^2\end{pmatrix}.
 \label{eq:activeblock}
\end{equation}
The imaginary block is identical. Each has eigenvalues \(0,-2\);
the zero direction is \((u,-b)\), consistent with
\eqref{eq:tangent}. The four real active off-diagonal directions
are tangents to the equality family and have zero Hessian.

Let \(\alpha\geq3\) be a spectator index.
For the real pair \((A_{1\alpha},A_{\alpha1})\),
the first-derivative vectors \(H_i e\) point only along
\(E_{\alpha2}\), with coefficients
\begin{equation}
 \frac{(1+q)b+2qu}{2},\qquad \frac{(1+q)b}{2}.
\end{equation}
The branch gap is
\begin{equation}
 cR-b^2=u^2(2c-q)>0.
\end{equation}
Using \((1+q)b+2qu=2uc\) in
\eqref{eq:lambdahess}--\eqref{eq:quotienthess} yields
\begin{equation}
 -\frac{a_1}{2}
 \begin{pmatrix}1&-1\\-1&1\end{pmatrix}.
 \label{eq:firstblock}
\end{equation}
For the real pair \((A_{2\alpha},A_{\alpha2})\), the
derivative vectors point only along \(E_{1\alpha}\), with
coefficients
\begin{equation}
 -\frac{(1+q)u}{2},\qquad
 -\frac{2b+(1+q)u}{2}.
\end{equation}
Now the gap and simplifying relation are
\begin{equation}
 cR-qu^2=b^2(2c-1)>0,\qquad
 2b+(1+q)u=2bc.
\end{equation}
The result is \eqref{eq:firstblock} with \(a_1\) replaced
by \(a_2\). For each imaginary pair the off-diagonal
signs reverse, so the nonzero eigenvalue is unchanged.
These zero directions are precisely the unitary variations
of the supporting two-plane.

For a direction \(D\) wholly within the spectator block,
\(DE_{12}=E_{12}D=0\), and \(H'_D e=0\).
Only the norm correction in \eqref{eq:quotienthess}
survives, giving \(-2c/R=-a_3\) for each real component.
There are \(2(n-2)^2\) such components.

There are no mixed Hessian blocks. Distinct spectator indices
couple to orthogonal matrix units; the two active-spectator
sectors couple respectively to \(E_{\alpha2}\) and
\(E_{1\alpha}\); active derivatives have only active support.
The direct-product cross terms in \eqref{eq:lambdahess} and
the norm cross terms vanish between these sectors as well.
This support argument holds for every finite \(n\).
It is not an inference from a finite list of computed dimensions.

The spectrum is therefore \eqref{eq:spectrum}. Since
\(c>\max(1,q)\), \(R>0\), \(2c-q>0\) and \(2c-1>0\),
all the displayed normal magnitudes are strictly positive.
The negative rank and nullity are
\begin{equation}
 2+4(n-2)+2(n-2)^2=2(n-1)^2,\qquad
 2n^2-2(n-1)^2=4n-2.
\end{equation}
Stationarity on \(\Mq\) puts its tangent space in the
Hessian kernel. Equality of dimensions identifies that
kernel exactly.

\section{Control of the full local remainder}
\begin{proof}[Completion of the proof of Theorem~\ref{thm:local}]
The simple largest eigenvalue of a real-analytic Hermitian
matrix family is real analytic locally. Since \(A_q\neq0\),
the quotient \(g_q\) has the same property.
In a sufficiently small tubular neighborhood write
\(A=p+v\), where \(p\in\Mq\) and \(v\) is
Frobenius-normal to \(\Mq\) at \(p\).
Then \(\dist(A,\Mq)=\norm{v}\).
Lemma~\ref{lem:family} gives
\(g_q(p)=c\) and \(dg_q(p)=0\).

At the base, the Hessian of \(h=c-g_q\) is positive
on the normal space, with smallest and largest eigenvalues
\(\mu,\Lambda\). Continuity of the Hessian and of the
normal bundle permits shrinking the neighborhood so that,
for every such normal segment and \(0\leq t\leq1\),
\begin{equation}
 \frac{\mu}{4}\norm{v}^2
 \leq D^2h(p+tv)[v,v]\leq
 2\Lambda\norm{v}^2.
 \label{eq:continuity}
\end{equation}
Here \(v\) is normal at \(p\); it need not be normal
to a different tangent plane at \(p+tv\).
Taylor's integral identity is exact:
\begin{equation}
 h(p+v)=\int_0^1(1-t)D^2h(p+tv)[v,v]\,dt.
\end{equation}
It proves \eqref{eq:deficit}. The radial derivative
\(dh(p+v)[v]\) is strictly positive whenever \(v\neq0\),
by the corresponding integral without the factor \(1-t\).
There are therefore no other nearby critical points.
The equality manifold is exactly the local maximum set,
and the Hessian is nondegenerate in its normal directions:
this is the Morse--Bott assertion.
Finally, \eqref{eq:g} proves \eqref{eq:localf}
for every \(B\), with \(B=0\) trivial.
\end{proof}

The argument controls the entire analytic remainder.
A negative quadratic truncation alone would not establish
a bound along the Hessian kernel; the exact equality
manifold is essential. The theorem also excludes a sequence
of violating pairs whose first operators approach \(A_q\)
at fixed \(n,q\), even if the second operators vary arbitrarily.

\section{Conditional-operation observable}
For nonzero matrices \(A,B\), choose positive scalars
\(\alpha=1/(2\norm{A})\) and \(\beta=1/(2\norm{B})\).
The operators \(K_A=\alpha A\) and \(K_B=\beta B\)
are interior contractions, since their operator norms
are at most \(1/2\).
Each is a legal success Kraus operator; adjoining
\((I-K^\dagger K)^{1/2}\) completes its two-outcome
trace-preserving instrument.

Use the normalized maximally entangled input
\begin{equation}
 |\Phi_n\rangle=\frac1{\sqrt n}\sum_{j=1}^n|j\rangle|j\rangle.
\end{equation}
With \(Y=K_BK_A\), \(X=K_AK_B\), define the successful
ordered amplitudes
\(|\psi_Y\rangle=(Y\otimes I)|\Phi_n\rangle\) and
\(|\psi_X\rangle=(X\otimes I)|\Phi_n\rangle\).
The exact trace identity
\begin{equation}
 \ip{\psi_Y}{\psi_X}=\frac{\Tr(Y^\dagger X)}n
 \label{eq:trace-map}
\end{equation}
fixes the normalization.
Assume coherent control of these two orders and retain the
joint-success branch. Its unnormalized state is
\begin{equation}
 |\Omega\rangle=\frac{|0\rangle|\psi_Y\rangle+
                          |1\rangle|\psi_X\rangle}{\sqrt2}.
\end{equation}
The reduced, subnormalized control state \(\rho\) obeys
\begin{equation}
 \begin{gathered}
 \rho_{00}=\frac{\norm{Y}^2}{2n},\quad
 \rho_{11}=\frac{\norm{X}^2}{2n},\quad
 \Tr(\sigma_x\rho)=\frac{\operatorname{Re}\Tr(Y^\dagger X)}n .
 \end{gathered}
\end{equation}
Thus the Hermitian control observable
\begin{equation}
 O_q=2|0\rangle\langle0|+2q|1\rangle\langle1|
                         -(1+q)\sigma_x
\end{equation}
has success-weighted response
\begin{equation}
 \mathcal R_q=\Tr(O_q\rho)
   =\frac{f(K_A,K_B;q)}n.
 \label{eq:observable}
\end{equation}
Failure events contribute zero to this response.
An expectation in the normalized postselected state would
instead be \(\mathcal R_q/\Tr\rho\), when the success
probability is nonzero, and is a different observable.
Negative \(q\) is a signed measurement weight, not a
negative probability or transition rate.

For \(A\) in Theorem~\ref{thm:local}'s neighborhood,
homogeneity transfers its bound to the physical operators:
\begin{equation}
 \mathcal R_q\leq
 \frac{\norm{K_A}^2\norm{K_B}^2}{n}
 \left[c-\frac{\mu}{8}\dist(A,\Mq)^2\right].
 \label{eq:physicalbound}
\end{equation}
The distance in this expression is measured in the
original representative \(A\), before contraction scaling.
The normalized response ratio is unchanged by that scaling.
At the saturation witness, an infinitesimal perturbation in
a unit Frobenius normal eigenvector with Hessian eigenvalue
\(-a\) has optimized deficit
\(c-g_q(A_q+tD)=a t^2/2+O(t^3)\).
Equation~\eqref{eq:spectrum} therefore gives exact normal
susceptibilities, including both kinds of couplings to extra levels.
No particular implementation, success efficiency or experimental
error radius is predicted.

\section{Comparison, limitations and reproducibility}
The established two-dimensional bound supplies the active
equality family. The advance here is its stability in an
open, full-dimensional neighborhood in \(M_n(\mathbb C)\),
the geometry of all local maximizing first operators, and
the complete normal response spectrum.
The auxiliary two-dimensional equality identity in
Appendix~\ref{app:equality} is an extraction of the known
two-dimensional inequality and is not a separate novelty claim.

Known bounds for \(\norm{AB-qBA}^2\)
\citep{ckos2022norm,nobori2026} concern a different response.
For example, the immediate combination of the ordinary
commutator bound with the general \(q\)-norm bound by
Cauchy--Schwarz gives \(\sqrt3\) at \(q=1/2\),
whereas the sharp response constant is
\((3+\sqrt{10})/4<\sqrt3\).
This simple implication is insufficient; we do not claim
that every possible transformation of prior results has
been exhausted. Likewise, the proved universal relaxation
bound \citep{muratore2024} is not this unaveraged all-\(q\)
response problem. A bounded primary-literature comparison
identified no matching local theorem; it is not exhaustive
priority certification.

The result is finite-dimensional and dimensionless.
There is no continuum limit, boundary condition, hidden
perturbative approximation or finite-size extrapolation.
The excluded values \(q=-1,1\) have known global bounds
but do not satisfy the nondegenerate geometry used here.
No explicit radius is obtained, and the theorem does not
settle \eqref{eq:global} for general finite displacements.
Determining an effective radius or extending control
beyond the local equality family would require further work.

\subsection{Symbolic checks and use of AI tools}
The analytic argument is accompanied by exact SymPy~1.14.0
checks in Python~3.12. The supplied programs construct
\eqref{eq:H} from multiplication operators, verify the
Schur equality pair, and reduce rational polynomial
residuals modulo \(u^2-b^2-2ub+2=0\).
Separate programs check the determinant identity in the
appendix and complete real-coordinate Hessians at
\(q=1/2\) in dimensions three and four.
The every-dimensional conclusion follows from the
matrix-unit support argument above, not from those finite
computations. The checks use exact identities, not
floating-point fitting or an enclosure inferred from a scan.
Their environment and regeneration instructions accompany
the local draft.

OpenAI Codex was used in literature screening, proposing
and developing arguments, implementing and executing
symbolic calculations, checking the mathematics and
operator interpretation, and drafting this manuscript.
The automated consistency checks do not constitute
external peer review or completed human scientific review.
Human author review and approval of this local draft remain pending.

\appendix
\section{Complete two-dimensional equality criterion at \texorpdfstring{\(q=1/2\)}{q=1/2}}
\label{app:equality}
\begin{proposition}
Let \(q=1/2\) and \(c=(3+\sqrt{10})/4\).
For every nonzero \(A\in M_2(\mathbb C)\),
\(g_{1/2}(A)=c\) if and only if
\begin{equation}
 \det A+\frac32(3+\sqrt{10})(\Tr A)^2=0.
 \label{eq:auxcriterion}
\end{equation}
\end{proposition}
\begin{proof}
Put \(D=c\norm{A}^2I_4-H_{1/2}(A)\) and
\(\Phi(A)=\det A+\tfrac32(3+\sqrt{10})(\Tr A)^2\).
The known two-dimensional theorem makes \(D\) positive
semidefinite. Direct determinant expansion yields
\begin{equation}
 \det D=\frac{\norm{A}^4}{64}|\Phi(A)|^2.
 \label{eq:detidentity}
\end{equation}
To perform the expansion for a general matrix, use unitary
Schur form \(A=\left(\begin{smallmatrix}a&t\\0&d\end{smallmatrix}\right)\),
where \(a,d\in\mathbb C\) and \(t\geq0\).
In the column-vectorized basis the response matrix is
\begin{equation}
 \begin{pmatrix}
 t^2 &-\overline a t/4&t(d-3a/4)&-3t^2/4\\
 -at/4&h_{22}&0&td/4\\
 t(\overline d-3\overline a/4)&0&h_{33}&
                         t(\overline a/2-3\overline d/4)\\
 -3t^2/4&t\overline d/4&t(a/2-3d/4)&t^2/2
 \end{pmatrix},
 \label{eq:schurH}
\end{equation}
with
\begin{align*}
 h_{22}&=|a|^2+\tfrac12|d|^2+\tfrac32t^2
                -\tfrac32\operatorname{Re}(\overline a d),\\
 h_{33}&=|d|^2+\tfrac12|a|^2
                -\tfrac32\operatorname{Re}(\overline a d).
\end{align*}
Subtraction from \(c(|a|^2+|d|^2+t^2)I_4\)
and expansion, with \((\sqrt{10})^2=10\), gives
\eqref{eq:detidentity}. Both sides are invariant under
unitary conjugation, so Schur reduction proves it for
every complex \(A\). Positive semidefiniteness and
\(A\neq0\) imply that its vanishing is equivalent
to the largest response eigenvalue being
\(c\norm{A}^2\). This includes nonzero nilpotent matrices.
\end{proof}

\section*{Statements and Declarations}
\paragraph{Data and code availability.}
No experimental data are used. Exact verification programs
and reproduction instructions are included in the local
handoff materials. A public archival location has not yet
been assigned.
\paragraph{Funding and competing interests.}
These declarations require confirmation by the human author
before submission.
\paragraph{Authorship.}
The stated name, affiliation and correspondence details,
as well as responsibility for the final scientific content,
require final author confirmation. AI assistance is described
in the methods discussion above.

\bibliographystyle{unsrtnat}
\bibliography{manuscript/references}
\end{document}
